\section{Capabilities and Challenges}
\label{sec:capabilities_challenges}

\subsection{System Capabilities}

Despite the limitations identified above, the current system provides a functional end-to-end pipeline that serves as the foundation for the proposed improvements.
The following capabilities are already present and require no redesign:

\begin{itemize}
    \item \textbf{Multi-modal search}: The system supports both text-based and image-based queries.
    Text queries are encoded by the BLaIR encoder and retrieved via the FAISS flat index, while image queries are encoded by CLIP and retrieved from a separate CLIP HNSW index.
    The two modalities share the same result presentation layer.

    \item \textbf{Hybrid re-ranking}: BM25 results from Tantivy and dense results from FAISS are already fused via Reciprocal Rank Fusion ($k = 60$), providing complementary exact-match and semantic recall.

    \item \textbf{User interaction tracking}: The \texttt{/interact} route records all click and rating events to MongoDB with millisecond-precision timestamps, building a complete and queryable interaction history for every authenticated user.

    \item \textbf{Exact-recall vector search}: The flat FAISS index provides exact nearest-neighbour retrieval over the full 3-million-item catalogue, serving as a reliable evaluation baseline and as a fallback when latency is not a constraint.

    \item \textbf{Modular API structure}: The FastAPI router architecture cleanly separates search, recommendation, interaction, profile, and authentication concerns, making individual modules replaceable without disrupting others.

    \item \textbf{Containerised infrastructure}: MongoDB and Redis are deployed as Docker containers, ensuring consistent behaviour across development and deployment environments.
\end{itemize}

\subsection{Identified Challenges}
\label{sec:challenges}

Table~\ref{tab:challenges} summarises the principal limitations identified through system analysis, grouped by module and linked to the proposed solution in Chapter~4.

\begin{table}[H]
\centering
\caption{Summary of identified challenges by module.}
\label{tab:challenges}
\begin{tabular}{p{2.8cm} p{5.8cm} p{3.8cm}}
\hline
\textbf{Module} & \textbf{Challenge} & \textbf{Addressed in} \\
\hline
Data Storage      & 612 ms P50 FAISS search latency; no programmatic indexes on MongoDB collections, causing full collection scans on all profile and interaction queries & Section~\ref{sec:storage_improvements} \\
\hline
Data Visualisation & Synchronous rendering causes UI freezes on slow queries; no filtering or sorting controls; no multilingual input support & Section~\ref{sec:viz_improvements} \\
\hline
Recommendation    & GRU-SeqDQN HR@10 = 0.1031 (near random); no content feature integration; no behavioural graph signal; large action-space convergence failure & Section~\ref{sec:infra_improvements} \\
\hline
Search latency    & End-to-end warm-cache latency of 1{,}102 ms driven by slow FAISS search and unoptimised encoder & Section~\ref{sec:infra_improvements} \\
\hline
Multilingual      & BLaIR encoder achieves Precision@10 $\leq 0.05$ on Vietnamese queries; no translation stage & Section~\ref{sec:infra_improvements} \\
\hline
\end{tabular}
\end{table}

Three challenges deserve particular emphasis because they have the largest impact on user-facing quality.

\paragraph{Recommendation quality —}
The GRU-SeqDQN model's near-random performance (HR@10 = 0.1031) renders the recommendation feature effectively unusable.
A user who sees random recommendations quickly loses trust in the system and stops engaging with the personalisation features entirely.
This is the highest-priority challenge and motivates the complete replacement of the recommendation backend with the dual-mode Pipeline A + Pipeline B architecture.

\paragraph{Search latency —}
A warm-cache end-to-end latency of 1{,}102 ms is above the 200 ms threshold at which users typically perceive a response as instantaneous.
Reducing this to below 100 ms requires replacing the flat FAISS index with an HNSW approximate index and adding GPU-accelerated encoding.

\paragraph{Multilingual gap —}
The BLaIR encoder's failure on Vietnamese queries (Precision@10 $\leq 0.05$, compared to 0.55 for BGE-M3 on long queries) creates an inequitable search experience for Vietnamese-speaking users who represent a significant portion of the target audience.
Addressing this challenge requires replacing BLaIR with a multilingual encoder and adding a query translation stage.

